\chapter{Cerumen Removal}

\section*{What Is This Test or Procedure?}
Cerumen removal is the extraction of impacted earwax from the external auditory canal to relieve symptoms or improve visualization of the tympanic membrane. It is performed with instruments, suction, irrigation, or cerumenolytic agents.

\section*{Alternative Names}
Earwax removal; ear irrigation; cerumen disimpaction; aural toilet

\section*{Principle}
Impacted cerumen is softened or displaced so it can be extracted from the ear canal. Instrumentation uses curettes or loops under direct visualization, irrigation uses a controlled water stream to flush wax outward, and cerumenolytics chemically soften the wax before removal. The procedure restores canal patency and auditory conduction.

\section*{History}
Cerumen removal has been practiced since antiquity, with ear syringing documented for centuries and modern suction and micro-instrumentation developed in the 20th century.

\section*{Why This Test or Procedure Is Ordered}
Performed for symptoms of cerumen impaction such as hearing loss, fullness, tinnitus, itching, or dizziness, and to allow examination of the tympanic membrane when wax obscures the view. It is also done to permit hearing aid fitting and to monitor the ear after surgery.

\section*{Is This a Primary or Secondary Test or Procedure}
A primary therapeutic procedure for symptomatic cerumen impaction.

\section*{How This Test or Procedure Is Performed}
The canal is inspected with an otoscope under magnification. The clinician removes wax using a loop or curette, suction, or irrigation with body-temperature water after confirming an intact tympanic membrane. A cerumenolytic may be instilled several days beforehand to soften hard wax.

\section*{What Preparation Is Needed}
No routine preparation is needed. Patients are asked about otorrhea, otalgia, prior ear surgery, and hearing aid use. An intact tympanic membrane is confirmed before irrigation.

\section*{Contraindications and Precautions}
Irrigation is contraindicated when the tympanic membrane is perforated, when a cholesteatoma or myringotomy tube is present, and in patients with a history of otitis media, because of the risk of infection or vertigo. Instrumentation is performed cautiously in anticoagulated patients and in those with a narrow or tortuous canal.

\section*{Risks}
Canal abrasion, bleeding, otitis externa, tinnitus, and tympanic membrane perforation; irrigation may rarely cause tympanic perforation, acute otitis media, or caloric-induced vertigo.

\section*{What Happens After the Test or Procedure}
The ear is re-examined to confirm complete removal and to assess the tympanic membrane. Patients are advised to keep the ear dry and to avoid inserting cotton swabs, which can push wax deeper.

\section*{Understanding Your Results}
Complete removal with an intact tympanic membrane and improved hearing confirms success. Persisting symptoms or inability to remove wax safely may indicate referral to an ENT specialist.

\section*{Normal Results}
A clear external auditory canal with a visible, intact tympanic membrane and no residual obstructing wax.

\section*{Follow-up Tests and Procedures}
If hearing loss persists despite clearing, audiometry is performed to evaluate for other causes. Patients with recurrent impaction may be advised to use cerumenolytics or periodic cleaning, and referral is considered for suspected underlying canal pathology.

\section*{References}
\begin{enumerate}
  \item MedlinePlus. Ear wax removal. U.S. National Library of Medicine; 2025.
  \item Current Medical Diagnosis and Treatment. McGraw-Hill; 2025.
\end{enumerate}

\clearpage
